\section{Introduction}
\label{sec:intro}

Consider an engineer asking whether a device can return to service after a
failed calibration. A current operating procedure, a past incident report and
a manual for a similar device can all contain relevant-looking text. These sources serve different purposes: the procedure is the intended source for the decision,
whereas the report explains what happened in an earlier case. An agent must
identify the collection intended for the task, use the managed metadata to
find a document that appears relevant and read the passage that addresses the
question. HMRS helps make those steps visible; it does not determine whether a
source governs the decision, is authoritative or is up to date. Even if the
right document has been filed correctly, a stale index or an incomplete read
of the source can leave that passage out of the evidence collected for the
answer.

Dense retrieval ranks text by learned relevance to a query \citep{dpr}, and
collection selection and scoped search have long been studied in distributed
information retrieval \citep{callan1995searching,si2003redde}. A ranked list from an unscoped
index offers candidate passages, but it does not show whether the agent
searched the intended collection, whether the applicable document was
visible to search, or which source passages the agent read. Enterprise content management already maintains
membership and descriptive records \citep{paivarinta-ecm-2005,tyrvaainen-ecm-2006},
but those records are only useful to the agent if it can consult them when
choosing where and what to search. The contribution here is not a new
collection-selection score: it brings recorded collection assignments,
descriptions and checks of search visibility after changes into one
agent-facing retrieval workflow.

We present the Hierarchical Management Retrieval System (HMRS), a self-hosted
platform that uses the same document organisation for both maintenance and
retrieval. During ingestion, an agent consults collection descriptions and
uses the recovered content to select a destination and derive document
descriptions and tags. The catalogue and document listings then expose that
placement and metadata to the query agent. In catalogue-first access, the
agent selects a collection for the task, inspects its document listing
(the shelf), and opens selected documents or parts to obtain passages with
source references. An optional passage scorer can filter that evidence.
When the scope is known, a pinned search can start directly in the specified
collections. The hierarchy guides where to look; lexical and vector search
remain available to find candidate text.

\ifdefined\HMRSReviewLayout\Needspace{5\baselineskip}\fi
\noindent\textbf{Contributions.}
\begin{enumerate}
\item \textbf{A managed retrieval control plane.} HMRS maintains
collections, documents or parts, and indexed segments in a shared management
model. The ingestion workflow uses source content to assist placement and
description, and lifecycle checks distinguish stored state from what is visible
to search (Sections~\ref{sec:system}--\ref{sec:organization}). HMRS does not
claim a learned classifier for collection ownership; the platform records and
exposes the placement decision made by the ingestion workflow.
\item \textbf{An auditable, task-directed access path.} The retrieval
agent uses catalogue and shelf descriptions to choose a collection and its
documents before reading selected source text. That path complements pinned
search over the same managed collection state. A separate, optional evidence
gate can admit passages without treating its score as a document ranking
(Section~\ref{sec:retrieval}). In this paper, source references are retrieval
locators for inspecting returned text, not claims of authority or entailment.
\item \textbf{Evaluation of distinct failure points.} Lifecycle probes,
scripted routing tests, evidence-pack comparisons and a fixed-pool replay
separate index visibility, collection choice, target-text coverage and
document order. A logged case illustrates the catalogue-first path, while
a separate blind comparison describes answer quality under native deployment
conditions (Section~\ref{sec:eval}).
\end{enumerate}

This work links document management and retrieval and makes each stage
inspectable. The evaluation does not test for a causal advantage over a
tool-matched flat or faceted retriever, and it does not establish claim-level
evidence support. The native answer-quality comparison is a deployment case
study; its configurations did not use matched search scopes or source access.
We therefore examine management state, collection choice, target-text coverage,
candidate ordering and answer quality separately; operational traces are used
to assess what can be reported about search scope and completion.

\subsection{Research questions and reading guide}
\label{sec:rq}

\textbf{RQ1} asks whether HMRS can distinguish stored state, indexed state
and query visibility under the tested probes after acquisition and selected
lifecycle operations.
\textbf{RQ2} asks how the tested catalogue-card configurations affect
collection selection and which conclusions can be drawn from the routing result.
\textbf{RQ3} asks how often the evidence pack contains the literal target
string as search scope and evidence assembly vary.
\textbf{RQ4} asks whether a passage-verifier score is suitable for document
ordering on fixed candidate pools. \textbf{RQ5} examines what operational
traces, interruption reports and instructed abstention reveal about search
completion. \textbf{RQ6} describes the answer-quality ratings from the native deployments, which were not matched for a controlled end-to-end comparison.

Section~\ref{sec:related} relates HMRS to document management, collection
selection and structured retrieval. Sections~\ref{sec:system}--\ref{sec:implementation}
explain how the shared hierarchy is maintained and queried. Section~\ref{sec:eval}
reports the evaluation, followed by the discussion, limitations and conclusion
in Sections~\ref{sec:discussion}--\ref{sec:conclusion}.
These questions deliberately decompose the pipeline into managed state,
collection choice, evidence coverage, ordering, completion and answer quality.
No single end-to-end score is used to stand in for these distinct observations.
